% ===========================================================================
%  Chapter 2 -- Data and Methods
%  Figures for this chapter go in figures/chapter2/
% ===========================================================================
\chapter{Data and Methods}\label{ch:data}
\chapterfigs{chapter2}

\section{Sample selection}\label{sec:sample}

The sample (\cref{fig:sample}, \cref{tab:sample}) is defined by stellar mass and \SFR, estimated with \textsc{cigale} \citep{Boquien2019}.

\begin{figure}[tbp]
  \centering
  \includegraphics[width=0.65\textwidth]{sample_mock}
  \caption{Distribution of the (mock) sample in the \Mstar--\SFR plane.}
  \label{fig:sample}
\end{figure}

\section{Tables}\label{sec:tables}

Tables use the thesis style defined in \texttt{preamble.tex}. Three environments are available: \texttt{uatblr} (plain), \texttt{uatalltblr} (with table notes, used inside \texttt{table}) and \texttt{ualongtblr} (multi-page, see \cref{app:long}). They are built on \texttt{tabularray}, so every \texttt{tabularray} key still works: \texttt{column\{1\}=\{font=\textbackslash bfseries\}} for a bold first column, \texttt{row\{even\}=\{bg=tablegray\}} for zebra rows, \texttt{Q[si=\{...\}]} for decimal alignment, and \verb|\TblrNote| for notes. In \cref{tab:sample} the header has two rows (names and units) and the last row is a summary row. 

\begin{table}[tbp]
  \centering
  \begin{uatalltblr}[
      caption = {Properties of the example sample (mock values).},
      label   = {tab:sample},
      note{a} = {Stellar mass derived from SED fitting.},
      note{b} = {Specific SFR averaged over the last \qty{100}{\mega\year}.},
      remark{Notes} = {All values are for illustration only.},
    ]{
      width   = 0.9\linewidth,
      colspec = {X[l] Q[si={table-format=1.4}] Q[si={table-format=2.2}] Q[si={table-format=-2.2}] X[c]},
      row{1-2}   = {guard, font=\tabfont\bfseries},
      row{2}     = {font=\tabfont},
      column{1}  = {font=\tabfont\bfseries},
      hline{2}   = {0pt},
      hline{3,7} = {0.05em},
    }
    ID    & {$z$}  & {$\log\Mstar$\TblrNote{a}} & {$\log\sSFR$\TblrNote{b}} & Environment \\
          &        & {(\unit{\solarmass})}        & {(\unit{\per\year})}       &             \\
    G001  & 0.0123 & 10.35 & -9.80  & Cluster \\
    G002  & 0.0210 &  9.87 & -9.41  & Field   \\
    G003  & 0.0187 & 10.92 & -10.25 & Cluster \\
    G004  & 0.0305 & 10.10 & -9.55  & Group   \\
    Mean  & 0.0206 & 10.31 & -9.75  & ---     \\
  \end{uatalltblr}
\end{table}


\section{Equations}\label{sec:eqs}

Equations are numbered by chapter; give each one a label and refer to it with \verb|\cref{eq:sfms}| (\cref{sec:xref}). This is an example: the main sequence is commonly parametrized as
\begin{equation}\label{eq:sfms}
  \log\SFR = \alpha\,\log\Mstar + \beta,
\end{equation}
with slope $\alpha$ and intercept $\beta$ \citep{Speagle2014}. Specific \SFR is $\sSFR=\SFR/\Mstar$.
